\begin{definitionbox}{Definition --- Metric Sequential Convergence}
\begin{definition}[Metric Sequential Convergence]
\label{def:metric-sequential-convergence}
Let \((X,d)\) be a metric space.  A sequence \(\{x_n\}\) in \(X\) is
\emph{convergent in \(X\)} if there exists \(x_0\in X\) such that for every
\(\varepsilon>0\), there exists \(N\in\mathbb{N}\) such that
\[
  d(x_n,x_0)<\varepsilon
  \qquad\text{for all } n\geq N.
\]
In this case, \(\{x_n\}\) converges to \(x_0\), written
\[
  x_n\to x_0,
  \qquad
  x_0=\lim_{n\to\infty}x_n.
\]
\end{definition}
\end{definitionbox}

\begin{remark*}[Standard quantified statement]
\[
\begin{aligned}
x_n\to x_0
\quad\Longleftrightarrow\quad
&x_0\in X
\;\land\;
\forall\varepsilon>0\,\exists N\in\mathbb{N}\,\forall n\in\mathbb{N}{:}
\\
&n\geq N
\Longrightarrow
d(x_n,x_0)<\varepsilon.
\end{aligned}
\]
\end{remark*}

\begin{remark*}[Predicate reading]
\[
\operatorname{ConvergesTo}(x_n,x_0,(X,d))
\]
means that the sequence \(x_n\) eventually lies inside every metric tolerance
around \(x_0\).
\end{remark*}

\begin{remark*}[Interpretation]
A sequence converges to \(x_0\) when its tail eventually lies inside every
metric tolerance around \(x_0\).  The metric supplies the meaning of
``within \(\varepsilon\)'' in the ambient space \(X\).
\end{remark*}

\begin{remark*}[Historical note]
This convention follows Definition~2.20 in Kainth's
\emph{A Comprehensive Textbook on Metric Spaces}
\cite{KainthComprehensiveMetricSpaces2023}.
\end{remark*}

\begin{dependencies}
\begin{itemize}
  \item \hyperref[def:metric-space]{Definition: Metric Space}
\end{itemize}
\end{dependencies}

\begin{definitionbox}{Definition --- Metric Neighborhood}
\begin{definition}[Metric Neighborhood]
\label{def:metric-neighborhood}
Let \((X,d)\) be a metric space and let \(x\in X\).  A subset \(U\subseteq X\)
is a \emph{neighborhood of \(x\)} if there exists \(\delta>0\) such that
\[
  B_d(x;\delta)\subseteq U.
\]
\end{definition}
\end{definitionbox}

\begin{remark*}[Standard quantified statement]
\[
U\text{ is a neighborhood of }x
\quad\Longleftrightarrow\quad
U\subseteq X
\;\land\;
\exists\delta>0{:}\ B_d(x;\delta)\subseteq U.
\]
\end{remark*}

\begin{remark*}[Interpretation]
A metric neighborhood of \(x\) is any subset of \(X\) that contains at least
one open ball centered at \(x\).  The set \(U\) may be larger than that ball;
the essential condition is that \(x\) has some positive-radius room inside
\(U\).
\end{remark*}

\begin{remark*}[Historical note]
This convention follows Definition~2.21 in Kainth's
\emph{A Comprehensive Textbook on Metric Spaces}
\cite{KainthComprehensiveMetricSpaces2023}.
\end{remark*}

\begin{dependencies}
\begin{itemize}
  \item \hyperref[def:metric-space]{Definition: Metric Space}
  \item \hyperref[def:metric-open-ball]{Definition: Open Ball}
\end{itemize}
\end{dependencies}

\begin{definitionbox}{Definition --- Epsilon Neighborhood}
\begin{definition}[Epsilon Neighborhood]
\label{def:metric-epsilon-neighborhood}
Let \((X,d)\) be a metric space, let \(x\in X\), and let
\(\varepsilon>0\).  The \emph{\(\varepsilon\)-neighborhood of \(x\)} is the
open ball
\[
  B_d(x;\varepsilon)
  =
  \{y\in X:d(y,x)<\varepsilon\}.
\]
\end{definition}
\end{definitionbox}

\begin{remark*}[Standard quantified statement]
\[
y\in B_d(x;\varepsilon)
\quad\Longleftrightarrow\quad
y\in X\;\land\; d(y,x)<\varepsilon.
\]
\end{remark*}

\begin{remark*}[Interpretation]
An \(\varepsilon\)-neighborhood is the concrete neighborhood obtained by
choosing a particular positive radius \(\varepsilon\).  It consists exactly of
the points whose distance from \(x\) is less than that tolerance.
\end{remark*}

\begin{dependencies}
\begin{itemize}
  \item \hyperref[def:metric-space]{Definition: Metric Space}
  \item \hyperref[def:metric-open-ball]{Definition: Open Ball}
  \item \hyperref[def:metric-neighborhood]{Definition: Metric Neighborhood}
\end{itemize}
\end{dependencies}

\begin{theorembox}{Theorem}
\begin{theorem}[Metric Neighborhood Criterion for Sequential Convergence]
\label{thm:metric-neighborhood-criterion-sequential-convergence}
Let \((X,d)\) be a metric space, let \(\{x_n\}\) be a sequence in \(X\), and
let \(x\in X\).  Then
\[
  x_n\to x
\]
if and only if every neighborhood of \(x\) contains all but finitely many terms
of \(\{x_n\}\).

\smallskip
\noindent
\hyperref[prf:metric-neighborhood-criterion-sequential-convergence]{\textit{Go to proof.}}
\end{theorem}
\end{theorembox}

\begin{remark*}[Standard quantified statement]
\[
x_n\to x
\quad\Longleftrightarrow\quad
\forall U\subseteq X{:}\;
\bigl(U\text{ is a neighborhood of }x\bigr)
\Longrightarrow
\exists N\in\mathbb{N}\,\forall n\geq N{:}\ x_n\in U.
\]
\end{remark*}

\begin{remark*}[Predicate reading]
\[
\operatorname{IsNeighborhood}(U,x,(X,d))
\Longrightarrow
\exists N\in\mathbb{N}\,\forall n\geq N{:}\ x_n\in U
\]
means that a convergent sequence is eventually contained in every
neighborhood of its limit.
\end{remark*}

\begin{remark*}[Interpretation]
Sequential convergence can be tested by neighborhoods instead of numerical
\(\varepsilon\)-bounds.  A convergent sequence eventually remains inside every
set that contains some open ball around the limit.
\end{remark*}

\begin{remark*}[Historical note]
This is the neighborhood characterization stated immediately after
Definition~2.21 in Kainth's \emph{A Comprehensive Textbook on Metric Spaces}
\cite{KainthComprehensiveMetricSpaces2023}.
\end{remark*}

\begin{dependencies}
\begin{itemize}
  \item \hyperref[def:metric-sequential-convergence]{Definition: Metric Sequential Convergence}
  \item \hyperref[def:metric-neighborhood]{Definition: Metric Neighborhood}
\end{itemize}
\end{dependencies}

\begin{definitionbox}{Definition --- Metric Cauchy Sequence}
\begin{definition}[Metric Cauchy Sequence]
\label{def:metric-cauchy-sequence}
Let \((X,d)\) be a metric space.  A sequence \(\{x_n\}\) in \(X\) is
\emph{Cauchy} if for every \(\varepsilon>0\), there exists \(N\in\mathbb{N}\)
such that
\[
  d(x_n,x_m)<\varepsilon
  \qquad\text{for all } n,m\geq N.
\]
\end{definition}
\end{definitionbox}

\begin{remark*}[Standard quantified statement]
\[
\{x_n\}\text{ is Cauchy}
\quad\Longleftrightarrow\quad
\forall\varepsilon>0\,\exists N\in\mathbb{N}\,\forall n,m\in\mathbb{N}{:}\;
n,m\geq N
\Longrightarrow
d(x_n,x_m)<\varepsilon.
\]
\end{remark*}

\begin{remark*}[Predicate reading]
\[
\operatorname{IsCauchy}(x_n,(X,d))
\]
means that the sequence \(x_n\) is eventually internally stable in the metric
space \((X,d)\).
\end{remark*}

\begin{remark*}[Interpretation]
A Cauchy sequence is eventually internally stable: sufficiently late terms are
close to one another, regardless of whether a limit point has already been
identified in \(X\).
\end{remark*}

\begin{remark*}[Historical note]
This convention follows Definition~2.22 in Kainth's
\emph{A Comprehensive Textbook on Metric Spaces}
\cite{KainthComprehensiveMetricSpaces2023}.
\end{remark*}

\begin{dependencies}
\begin{itemize}
  \item \hyperref[def:metric-space]{Definition: Metric Space}
\end{itemize}
\end{dependencies}

\begin{theorembox}{Theorem}
\begin{theorem}[Metric Convergent Sequences Have Unique Limits]
\label{thm:metric-convergent-sequences-unique-limits}
In metric spaces, convergent sequences have unique limits.

\smallskip
\noindent
\hyperref[prf:metric-convergent-sequences-unique-limits]{\textit{Go to proof.}}
\end{theorem}
\end{theorembox}

\begin{remark*}[Standard quantified statement]
\[
\forall \{x_n\}\subseteq X\,\forall x,y\in X{:}\quad
\bigl(x_n\to x\ \land\ x_n\to y\bigr)
\Longrightarrow
x=y.
\]
\end{remark*}

\begin{remark*}[Predicate reading]
\[
\operatorname{ConvergesTo}(x_n,x,(X,d))
\land
\operatorname{ConvergesTo}(x_n,y,(X,d))
\Longrightarrow x=y
\]
means that a metric sequence has at most one limit.
\end{remark*}

\begin{remark*}[Interpretation]
In a metric space, a convergent sequence cannot settle at two different
points.  The positive definiteness of the metric turns zero distance between
candidate limits into equality of the points.
\end{remark*}

\begin{remark*}[Historical note]
This result corresponds to Theorem~2.23 in Kainth's
\emph{A Comprehensive Textbook on Metric Spaces}
\cite{KainthComprehensiveMetricSpaces2023}.
\end{remark*}

\begin{dependencies}
\begin{itemize}
  \item \hyperref[def:metric-space]{Definition: Metric Space}
  \item \hyperref[def:metric-sequential-convergence]{Definition: Metric Sequential Convergence}
\end{itemize}
\end{dependencies}

\begin{theorembox}{Theorem}
\begin{theorem}[Metric Cauchy Sequence with Convergent Subsequence]
\label{thm:metric-cauchy-sequence-convergent-subsequence}
Let \(\{x_n\}\) be a Cauchy sequence in a metric space \((X,d)\), let
\(x\in X\), and let \(\{x_{n_k}\}\) be a subsequence of \(\{x_n\}\) such that
\[
  \lim_{k\to\infty}x_{n_k}=x.
\]
Then \(x_n\to x\).

\smallskip
\noindent
\hyperref[prf:metric-cauchy-sequence-convergent-subsequence]{\textit{Go to proof.}}
\end{theorem}
\end{theorembox}

\begin{remark*}[Standard quantified statement]
\[
\begin{aligned}
\forall \{x_n\}\subseteq X\,\forall x\in X{:}\quad
&\Bigl(
\{x_n\}\text{ is Cauchy}
\\
&{}\land
\exists\{x_{n_k}\}\text{ subsequence of }\{x_n\}{:}\ x_{n_k}\to x
\Bigr)
\Longrightarrow
x_n\to x.
\end{aligned}
\]
\end{remark*}

\begin{remark*}[Predicate reading]
\[
\operatorname{IsCauchy}(x_n,(X,d))
\land
\operatorname{IsSubsequence}(x_{n_k},x_n)
\land
\operatorname{ConvergesTo}(x_{n_k},x,(X,d))
\Longrightarrow
\operatorname{ConvergesTo}(x_n,x,(X,d))
\]
means that a Cauchy sequence follows any convergent subsequence to the same
limit.
\end{remark*}

\begin{remark*}[Interpretation]
For a Cauchy sequence, one convergent subsequence determines the behavior of
the whole sequence.  The Cauchy condition forces the full tail to follow any
subsequence tail that has already converged.
\end{remark*}

\begin{remark*}[Historical note]
This result corresponds to Theorem~2.24 in Kainth's
\emph{A Comprehensive Textbook on Metric Spaces}
\cite{KainthComprehensiveMetricSpaces2023}.
\end{remark*}

\begin{dependencies}
\begin{itemize}
  \item \hyperref[def:metric-space]{Definition: Metric Space}
  \item \hyperref[def:metric-cauchy-sequence]{Definition: Metric Cauchy Sequence}
  \item \hyperref[def:metric-sequential-convergence]{Definition: Metric Sequential Convergence}
\end{itemize}
\end{dependencies}

\begin{theorembox}{Theorem}
\begin{theorem}[Coordinatewise Sequences in \(\mathbb{R}^m\)]
\label{thm:metric-euclidean-coordinatewise-sequences}
Let \(\{x_n\}\) be a sequence in \(\mathbb{R}^m\), and let
\(x_0\in\mathbb{R}^m\).  Write
\[
  x_n=\bigl(x_n^{(1)},\ldots,x_n^{(m)}\bigr)
  \qquad\text{for all } n\in\mathbb{N}\cup\{0\}.
\]
Then:
\begin{enumerate}
  \item \(x_n\to x_0\) if and only if
        \(x_n^{(j)}\to x_0^{(j)}\) for every \(j=1,\ldots,m\).
  \item \(\{x_n\}\) is Cauchy if and only if
        \(\{x_n^{(j)}\}\) is Cauchy for every \(j=1,\ldots,m\).
\end{enumerate}

\smallskip
\noindent
\hyperref[prf:metric-euclidean-coordinatewise-sequences]{\textit{Go to proof.}}
\end{theorem}
\end{theorembox}

\begin{remark*}[Standard quantified statement]
\[
\begin{aligned}
x_n\to x_0
&\Longleftrightarrow
\forall j\in\{1,\ldots,m\}{:}\; x_n^{(j)}\to x_0^{(j)},\\
\{x_n\}\text{ is Cauchy}
&\Longleftrightarrow
\forall j\in\{1,\ldots,m\}{:}\; \{x_n^{(j)}\}\text{ is Cauchy}.
\end{aligned}
\]
\end{remark*}

\begin{remark*}[Predicate reading]
\[
\operatorname{ConvergesTo}(x_n,x_0,\mathbb{R}^m)
\Longleftrightarrow
\forall j\in\{1,\ldots,m\}{:}\;
\operatorname{ConvergesTo}(x_n^{(j)},x_0^{(j)},\mathbb{R})
\]
and the analogous Cauchy statement mean that finite-dimensional Euclidean
sequence behavior is coordinatewise.
\end{remark*}

\begin{remark*}[Interpretation]
In finite-dimensional Euclidean space, convergence and Cauchy behavior can be
checked one coordinate at a time.  The Euclidean metric packages the finitely
many coordinate estimates into one distance estimate.
\end{remark*}

\begin{remark*}[Historical note]
This result corresponds to Theorem~2.25 in Kainth's
\emph{A Comprehensive Textbook on Metric Spaces}
\cite{KainthComprehensiveMetricSpaces2023}.
\end{remark*}

\begin{dependencies}
\begin{itemize}
  \item \hyperref[def:metric-sequential-convergence]{Definition: Metric Sequential Convergence}
  \item \hyperref[def:metric-cauchy-sequence]{Definition: Metric Cauchy Sequence}
  \item \hyperref[def:euclidean-distance-rn]{Definition: Euclidean Distance on \(\mathbb{R}^n\)}
  \item \hyperref[cor:euclidean-metric]{Corollary: Euclidean Metric}
\end{itemize}
\end{dependencies}

\begin{theorembox}{Theorem}
\begin{theorem}[Euclidean Cauchy Sequences Converge]
\label{thm:metric-euclidean-cauchy-sequences-converge}
Every Cauchy sequence in \(\mathbb{R}^m\) is convergent in \(\mathbb{R}^m\).

\smallskip
\noindent
\hyperref[prf:metric-euclidean-cauchy-sequences-converge]{\textit{Go to proof.}}
\end{theorem}
\end{theorembox}

\begin{remark*}[Standard quantified statement]
\[
\forall \{x_n\}\subseteq\mathbb{R}^m{:}\quad
\{x_n\}\text{ is Cauchy}
\Longrightarrow
\exists x\in\mathbb{R}^m{:}\; x_n\to x.
\]
\end{remark*}

\begin{remark*}[Predicate reading]
\[
\operatorname{IsCauchy}(x_n,\mathbb{R}^m)
\Longrightarrow
\exists x\in\mathbb{R}^m{:}\;
\operatorname{ConvergesTo}(x_n,x,\mathbb{R}^m)
\]
means that Euclidean space contains limits of its Cauchy sequences.
\end{remark*}

\begin{remark*}[Interpretation]
\(\mathbb{R}^m\) contains the limits of all of its Cauchy sequences.  This is
the finite-dimensional Euclidean completeness property.
\end{remark*}

\begin{remark*}[Historical note]
This result corresponds to Theorem~2.26 in Kainth's
\emph{A Comprehensive Textbook on Metric Spaces}
\cite{KainthComprehensiveMetricSpaces2023}.
\end{remark*}

\begin{dependencies}
\begin{itemize}
  \item \hyperref[def:metric-cauchy-sequence]{Definition: Metric Cauchy Sequence}
  \item \hyperref[def:metric-sequential-convergence]{Definition: Metric Sequential Convergence}
  \item \hyperref[thm:metric-euclidean-coordinatewise-sequences]{Theorem: Coordinatewise Sequences in \(\mathbb{R}^m\)}
\end{itemize}
\end{dependencies}

\begin{theorembox}{Theorem}
\begin{theorem}[Bolzano--Weierstrass in \(\mathbb{R}^m\)]
\label{thm:metric-euclidean-bolzano-weierstrass}
Every bounded sequence in \(\mathbb{R}^m\) contains a subsequence that
converges in \(\mathbb{R}^m\).

\smallskip
\noindent
\hyperref[prf:metric-euclidean-bolzano-weierstrass]{\textit{Go to proof.}}
\end{theorem}
\end{theorembox}

\begin{remark*}[Standard quantified statement]
\[
\forall \{x_n\}\subseteq\mathbb{R}^m{:}\quad
\{x_n\}\text{ is bounded}
\Longrightarrow
\exists\{x_{n_k}\}\text{ subsequence }\exists x\in\mathbb{R}^m{:}\;
x_{n_k}\to x.
\]
\end{remark*}

\begin{remark*}[Predicate reading]
\[
\operatorname{Bounded}(x_n,\mathbb{R}^m)
\Longrightarrow
\operatorname{HasConvergentSubsequence}(x_n,\mathbb{R}^m)
\]
means that every bounded Euclidean sequence has a convergent subsequence.
\end{remark*}

\begin{remark*}[Interpretation]
Boundedness in finite-dimensional Euclidean space is strong enough to force
some subsequence to settle to a Euclidean limit.
\end{remark*}

\begin{remark*}[Historical note]
This result corresponds to Theorem~2.27 in Kainth's
\emph{A Comprehensive Textbook on Metric Spaces}
\cite{KainthComprehensiveMetricSpaces2023}.
\end{remark*}

\begin{dependencies}
\begin{itemize}
  \item \hyperref[def:metric-sequential-convergence]{Definition: Metric Sequential Convergence}
  \item \hyperref[def:euclidean-distance-rn]{Definition: Euclidean Distance on \(\mathbb{R}^n\)}
  \item \hyperref[cor:euclidean-metric]{Corollary: Euclidean Metric}
\end{itemize}
\end{dependencies}
